% Compile with: latexmk -pdf lean-search-bench.tex   (pdfLaTeX + BibTeX)
% Format follows the Lean Census and Fuse papers.
% figures/forms.pdf and scores.json are produced by figures.py.
\documentclass{article}
\PassOptionsToPackage{numbers,compress}{natbib}
\usepackage[preprint]{neurips_2025}
\usepackage[utf8]{inputenc}
\usepackage[T1]{fontenc}
\usepackage{microtype}
\usepackage{amsmath,amssymb}
\usepackage{booktabs}
\usepackage{tabularx}
\usepackage{enumitem}
\usepackage{xcolor}
\usepackage{colortbl}
\usepackage{graphicx}
\usepackage{float}

% Numina palette (shared with the Fuse and Lean Census papers).
\definecolor{nAccent}{HTML}{BB480B}
\definecolor{nInk}{HTML}{222222}
\definecolor{nMuted}{HTML}{6B5B4A}
\definecolor{nBorder}{HTML}{D6CFC0}
\definecolor{nBorderLight}{HTML}{EDE6DA}

\usepackage[font=small,labelfont={bf,color=nAccent},labelsep=period]{caption}
\captionsetup[table]{skip=8pt}
\usepackage[colorlinks=true,linkcolor=nAccent,citecolor=nAccent,urlcolor=nAccent]{hyperref}
\arrayrulecolor{nBorder}

\newcommand{\thead}[1]{\textbf{#1}}
\newcommand{\rowrule}{\arrayrulecolor{nBorderLight}\specialrule{0.4pt}{3pt}{3pt}\arrayrulecolor{nBorder}}
\renewcommand{\bottomfraction}{0.6}
\renewcommand{\textfraction}{0.1}
\renewcommand{\topfraction}{0.9}
\renewcommand{\floatpagefraction}{0.8}
\setlist{itemsep=2pt,topsep=3pt}
\raggedbottom
\newcommand{\code}[1]{\texttt{\small #1}}
\newcommand{\todo}[1]{\textcolor{red}{[TODO: #1]}}

\title{A Balanced Benchmark for Lean 4 Declaration Search}
\author{%
  Justin Asher$^{1}$\hspace{1.4em}Pietro Pellegrino$^{1,2}$ \\[1.2ex]
  \small\normalfont $^{1}$Project Numina \\[0.3ex]
  \small\normalfont $^{2}$University of Naples Federico II \\[0.8ex]
  \small\normalfont \texttt{justin@projectnumina.ai}\hspace{1.4em}\texttt{pietro@projectnumina.ai}
}

\begin{document}
\maketitle

\begin{abstract}
Search engines for Lean 4 declarations are evaluated on query sets built by
the authors of the engine under test, in the styles that engine handles. We
present a benchmark of 3,200 queries for search over Mathlib, in four buckets
of 800 that are weighted equally: names and keywords, natural language, formal
statements and proof states, and questions derived from Lean Zulip threads.
The forms and lengths of the queries follow aggregate statistics of real
searches. Every answer is a Mathlib declaration that all engines index, so the
score measures search and not coverage. Answers are completed by pooling the
results of four engines and grading them blind. On the test split, Lean Finder
and LeanSearch find an answer in the first ten results for 92\% and 91\% of
queries, and LeanExplore, our own engine, for 70\%. We release the queries,
labels, engine results, and code at
\url{https://github.com/lean-explore/lean-explore-bench}.
\end{abstract}

%=====================================================================
\section{Introduction}

A person or an AI agent writing Lean~\cite{moura2021lean4} spends much of the
time looking for declarations that already exist. Mathlib~\cite{mathlib2019}
has several hundred thousand of them, and several engines now search it:
Loogle~\cite{breitner2023loogle} by name and pattern, and
LeanSearch~\cite{gao2024leansearch,gao2026leansearchv2},
LeanExplore~\cite{asher2025leanexplore}, and Lean
Finder~\cite{lu2025leanfinder} by meaning. LeanExplore alone received about
50,000 searches in six weeks of autumn 2026, three quarters of them from
programs and not from people at a keyboard.

There is no shared way to compare the engines. Each was evaluated by its
authors on a query set of their own design, in the styles of query that the
engine was built for. A benchmark for this task must settle three questions,
and each has an answer that favours some engine.

\paragraph{Which queries.} Real searches take many forms, from a bare
declaration name to a pasted proof state. An engine may be excellent on one
form and poor on another, so the mix of forms decides the ranking. Weighting
the forms by the traffic of one engine favours that engine, because people
learn what an engine handles and send it that.

\paragraph{Which declarations.} The engines index different libraries at
different dates. An engine that does not index a declaration cannot return it,
and a benchmark that asks for it measures coverage. Coverage is a fact about a
deployment that can be read from its index. How well an engine searches cannot.

\paragraph{Which answers.} A query written from one declaration is usually
labelled with that declaration alone, although Mathlib often has several that
answer it. An engine that returns one of the others is marked wrong.

We make one choice for each. The benchmark has four buckets, each a kind of
search that exists apart from any engine, and its score is the plain mean of
the four. Every answer is a Mathlib declaration that was already present in
January 2026, which all engines index. Answers are completed by pooling: the
results that four engines returned were graded without knowledge of engine or
rank, and those that answer the query were added as labels.

On this benchmark Lean Finder and LeanSearch are close to each other and well
ahead of LeanExplore, which we develop. LeanExplore is nearly level with them
on names and keywords and far behind on formal queries and on questions from
real discussions. We report this as we found it. A benchmark is credible
only if the engine of its authors can lose on it.

%=====================================================================
\section{What people search for}
\label{sec:searches}

We sorted searches into eight forms by surface features alone: which tokens
look like Lean names, how many function words such as ``of'' or ``the'' a
query has, and how much Lean syntax. The subject of a query plays no part.
Table~\ref{tab:forms} defines the forms.

\begin{table}[b]
\centering\small
\caption{The eight forms of a query, with an example from the benchmark and
the share of LeanExplore searches of each form from 29 September to 11
November 2026.}
\label{tab:forms}
\begin{tabularx}{\linewidth}{@{}llXr@{}}
\toprule
\thead{Form} & \thead{Definition} & \thead{Example} & \thead{Searches} \\
\midrule
\code{name} & One Lean name & \code{Polynomial.Separable.of\_mul\_right} & 5.5\% \\
\rowrule
\code{names} & Several Lean names & \code{IsOpen Set.singleton pure\_le\_nhds} & 1.5\% \\
\rowrule
\code{names+keywords} & Names and other words & \code{Finsupp.prod} embedding domain support & 22.7\% \\
\rowrule
\code{keywords} & Words, no function words & combine family affine maps into product space & 26.2\% \\
\rowrule
\code{prose} & Natural language & tail of interleaving two infinite streams & 23.9\% \\
\rowrule
\code{prose+names} & Natural language with names & prenex conversion of a \code{BoundedFormula} is prenex & 17.0\% \\
\rowrule
\code{formal} & A Lean statement or pattern & \code{ManyOneEquiv u v $\to$ (ManyOneEquiv u w $\leftrightarrow$ ManyOneEquiv v w)} & 2.8\% \\
\rowrule
\code{goal} & A proof state & \code{$\vdash$ WithBot.succ x = $\bot$ $\leftrightarrow$ x = $\bot$} & 0.4\% \\
\bottomrule
\end{tabularx}
\end{table}

\paragraph{How the search log was used.} We applied the same rules to
LeanExplore's search log, about 50,000 searches of which 42,000 are distinct.
The log was used only to count: how searches divide over the forms, how long
they are, how many Lean names they hold, and how often they are lower case or
phrased as questions. No logged search is in the benchmark, and none was shown
to a language model. The analysis ran on the server that holds the log and
returned shares only, with counts under ten left out. Search logs identify
people more easily than they appear to~\cite{barbaro2006aol}, and the counts
are all the benchmark needs.

\paragraph{What the counts show.} Half of all searches are names and keywords
with no sentence structure, and four in ten are natural language. Lean
statements and proof states are 3\%. Searches are short, with a median of one
word for a name, five or six for keyword forms, and ten to twelve for the
rest. The source matters. The programmatic interface accounts for 77\% of
searches and leans towards names with keywords (26\% against 10\% on the
website), while the website leans towards prose (46\% against 18\%).

These shares describe one engine. LeanExplore searches over names and
natural-language descriptions and is weak on formal input, as
Section~\ref{sec:results} shows, so its users have little reason to paste a
proof state into it. The 3\% is a fact about LeanExplore's traffic and not
about how often people want to search by statement. This is why the benchmark
does not use the shares as weights.

%=====================================================================
\section{Design}
\label{sec:design}

\paragraph{Four buckets with equal weight.} The benchmark groups the forms
by what the searcher is doing (Table~\ref{tab:buckets}). Someone who types
names and keywords half-knows the declaration. Someone who writes a sentence
knows the mathematics and not the name. Someone who pastes a statement or a
goal wants the lemma that proves it. The fourth bucket holds questions that
people asked each other. Each bucket has 800 queries, and \textbf{the
benchmark score is the plain mean of the four bucket scores}. Within a bucket
the forms are in the proportions of Table~\ref{tab:buckets}, which follow the
real shares loosely and in round numbers. A fifth of each bucket is a
development split of 640 queries, and the remaining 2,560 are the test split.

\begin{table}[t]
\centering\small
\caption{The four buckets. Median words and answers per query are over all
3,200 queries. An answer is a declaration graded 2 or 3.}
\label{tab:buckets}
\begin{tabular}{@{}lrlrr@{}}
\toprule
\thead{Bucket} & \thead{Queries} & \thead{Forms} & \thead{Median words} & \thead{Answers} \\
\midrule
Names and keywords & 800 & 160 \code{name}, 80 \code{names}, & 1 to 7 & 1.2 to 1.7 \\
 & & 240 \code{names+keywords}, 320 \code{keywords} & & \\
\rowrule
Natural language & 800 & 480 \code{prose}, 320 \code{prose+names} & 12 to 14 & 1.3 to 1.5 \\
\rowrule
Formal & 800 & 480 \code{formal}, 320 \code{goal} & 14 to 19 & 1.1 to 1.2 \\
\rowrule
Real questions & 800 & from Lean Zulip threads & 15 & 2.2 \\
\bottomrule
\end{tabular}
\end{table}

\paragraph{Mathlib only, and only established declarations.} Every engine
indexes Mathlib, and no other library is indexed by all of them, so every
answer is a Mathlib declaration. The engines also index Mathlib at different
dates. In an earlier build of the benchmark, on queries written from
declarations added to Mathlib after January 2026, LeanExplore found the answer
in the first ten results for 80\% of queries and LeanSearch and Lean Finder
for none, where on older declarations they found 93\% and 95\%. A declaration
is therefore an answer only if it has the same name and module in a January
2026 snapshot. Of the Mathlib declarations in our July 2026 corpus, 88\% do.

\paragraph{Splits.} Queries that are not independent share a group, and a
group is never divided between development and test. Two queries from one
discussion thread form a group, as do a query and its rewrite.

%=====================================================================
\section{Construction}
\label{sec:construction}

\subsection{Written queries}

The first three buckets are written from a declaration, which gives each
query one answer that is known to be correct. This is known-item
simulation~\cite{azzopardi2007simulated}, with a language model as the
writer~\cite{rahmani2024synthetic}.

\paragraph{Targets.} Targets are drawn at random from the established
declarations of Mathlib, with the draw flattened across its top-level areas so
that the largest do not dominate. Declarations that Lean generates, as opposed
to those an author wrote, are not eligible. The 2,400 written queries have 2,397
distinct targets.

\paragraph{Writing.} The writer sees the target's name, statement,
documentation, and a natural-language description, together with the form to
write and a style drawn from the real searches of that form: a word count,
whether the query is lower case, and whether it is a question. A
\code{name} query is produced by rule, as one of five variants: the full
name, the name without its leading namespaces, the name with dots and
underscores turned into spaces, the name with a typing error, or the name in
lower case. Lean statements and proof states
are written at the length of the statement. An earlier build held them to a
word count from real searches, which produced padded and truncated statements,
and we reviewed and rewrote those.

\paragraph{Checks.} A query passes rule checks before it is judged. Its form,
as decided by the rules of Section~\ref{sec:searches}, must be the form asked
for. It must not contain the target's name unless the form calls for names,
and it must be near the requested length.

\paragraph{Judge.} A model from another family than the writer grades each
query against its target, from 0 to 3. Grade 3 means that someone who searched
with this query would be satisfied with the declaration as it stands, and
grade 2 that they would have to adapt it in a clear way. The judge also says
whether the query is specific, meaning that only this declaration or
near-identical variants answer it. A written query is accepted at grade 3 and
specific.

\paragraph{Writers.} Of the 2,400 written queries, 160 names are made by rule
and 1,294 were written by Claude Opus. Another 946 were first written by
Claude Sonnet 5.5 and then reviewed by Claude Opus, which kept 624 and
rewrote 322. The faults that led to a rewrite were Lean identifiers that do
not exist, lemmas cited from the proof and not the statement, keyword
stuffing, and queries that a sibling lemma answers equally well. A rewrite
replaced a query only when the judge, \code{openai/gpt-6.1-terra}, accepted
it. Every record names its writer.

\subsection{Real questions}

The fourth bucket is not written from a declaration. Its queries come from
threads of the Lean Zulip chat in which someone asks whether a result exists
and someone answers with the name of a declaration.

\paragraph{From thread to query.} For each such thread a language model
writes the search that the person asking might have typed, using only what the
question says and nothing learned from the replies. The declarations named in the replies are
the proposed answers. The judge grades each proposed answer against the query,
and the query is kept when at least one has grade 2 or 3.

\paragraph{In its own words.} The benchmark contains no text of any post. A
query that repeats more than five consecutive words of a message in its
thread is rejected. For 515 queries the thread could not be retrieved, and
these have not been checked. \todo{check them before release}

\paragraph{Two collections.} The bucket combines two collections made
independently. Of 1,807 queries in the first, the judge accepted 767. Another
648 had no proposed answer that is an eligible declaration, 349 had every
proposed answer rejected, 34 contained the name of an answer, and 9 copied a
post. Of the 767, 177 have answers outside Mathlib or newer than January
2026, and 590 are in the benchmark. The second collection, from 1,650 threads,
supplies the other 210.

\subsection{Pooled answers}
\label{sec:pooling}

A written query starts with one answer, and a real question with the few
that its thread names. Pooling is the standard remedy for incomplete
labels~\cite{sparckjones1975pooling,voorhees2005trec}: the top results of
several systems are merged and judged, and what is found relevant is added.

We ran four engines on every query and pooled their first ten results. Each
engine's list was cut at its first labelled answer, and candidates outside
Mathlib or not established were removed. This left 20,141 candidates for
2,669 queries. Claude Opus graded each candidate on the scale above. It saw
the query, the answers already labelled, and the candidates in alphabetical
order with their source, and did not see which engine returned a candidate or
at what rank. Candidates graded 2 or 3 were added as labels.

Pooling added 1,551 answers to 984 queries, so that 1,089 of the 3,200
queries now have more than one answer. About two thirds of pooled queries
gained nothing. The labels hold 3,348 answers of grade 3, 1,668 of grade 2,
and 240 declarations of grade 1, which are related but do not count.

\paragraph{What the cut preserves.} Whether an engine has an answer among its
first $k$ results, and the rank of its first answer, depend only on its
results up to its first answer. If the first answer under complete labels is
the one already labelled, nothing above it was an answer. If it is an earlier
result, that result was in the pool. Success@$k$ and reciprocal rank are
therefore the same as if all ten results of every engine had been judged,
while about half as many candidates need a grade. nDCG is not preserved,
since it also counts answers below the first.

%=====================================================================
\section{Results}
\label{sec:results}

\paragraph{Setup.} We queried the hosted services of
LeanExplore~\cite{asher2025leanexplore},
LeanSearch~\cite{gao2026leansearchv2}, Lean Finder~\cite{lu2025leanfinder},
and Loogle~\cite{breitner2023loogle} from 11 to 13 November 2026 and kept the
first ten results of each. All 12,800 requests succeeded. A result is a hit
when its grade is 2 or 3. We report Success@10, the share of queries with a
hit in the first ten results, and Success@1, the share whose first result is a
hit.

\begin{table}[t]
\centering\small
\caption{Success@10 on the test split, with Success@1 in brackets. Each bucket
has 640 test queries. The interval is a 95\% bootstrap interval for the
benchmark score.}
\label{tab:results}
\begin{tabular}{@{}lrrrr@{}}
\toprule
\thead{Bucket} & \thead{LeanExplore} & \thead{LeanSearch} & \thead{Lean Finder} & \thead{Loogle} \\
\midrule
Names and keywords & 0.91 (0.52) & 0.93 (0.62) & 0.94 (0.75) & 0.07 (0.07) \\
\rowrule
Natural language & 0.86 (0.48) & 0.98 (0.72) & 0.98 (0.88) & 0.00 (0.00) \\
\rowrule
Formal & 0.54 (0.14) & 0.95 (0.68) & 0.96 (0.89) & 0.10 (0.10) \\
\rowrule
Real questions & 0.50 (0.27) & 0.77 (0.46) & 0.80 (0.56) & 0.00 (0.00) \\
\midrule
\textbf{Benchmark score} & \textbf{0.70 (0.35)} & \textbf{0.91 (0.62)} & \textbf{0.92 (0.77)} & \textbf{0.04 (0.04)} \\
Success@10 interval & 0.69 to 0.72 & 0.90 to 0.92 & 0.91 to 0.93 & 0.04 to 0.05 \\
Success@1 interval & 0.34 to 0.37 & 0.60 to 0.64 & 0.75 to 0.79 & 0.03 to 0.05 \\
\bottomrule
\end{tabular}
\end{table}

\paragraph{The ranking.} Table~\ref{tab:results} gives the scores. Lean
Finder and LeanSearch find an answer in the first ten results for 92\% and
91\% of queries, and their intervals overlap. They differ in where the answer
lands: Lean Finder puts an answer first for 77\% of queries and LeanSearch for
62\%. LeanExplore is behind both, at 70\% and 35\%. Loogle searches by name and
pattern and is not built for most of these queries, and we include it as a
reference point for the two forms that suit it.

\paragraph{By bucket.} LeanExplore is within three points of the others on
names and keywords at ten results, and twelve points behind on natural
language. On formal queries it finds an answer for 54\% against 95\% and 96\%,
and puts it first for 14\% against 68\% and 89\%. Real questions are the
hardest bucket for every engine. The best engine answers 80\% of them, and for
86 of the 640 no semantic engine returns an answer in ten results, against 19
of the 1,920 written test queries.

\paragraph{By form.} Figure~\ref{fig:forms} divides the buckets into forms.
The bare name is the one form on which LeanExplore leads at ten results (0.89
against 0.85 and 0.80), and Loogle answers 31\% of these, about the share that
are full names typed without error. On proof states LeanExplore puts the answer first for 6\% of
queries, against 65\% for LeanSearch and 90\% for Lean Finder. Proof states are one of
the inputs that Lean Finder was built and evaluated
for~\cite{lu2025leanfinder}, and
LeanExplore searches over names and natural-language descriptions and not over
statements.

\begin{figure}[t]
\centering
\includegraphics[width=\linewidth]{figures/forms.pdf}
\caption{Scores by form on the test split. \code{question} is the real
questions bucket. Loogle is left out.}
\label{fig:forms}
\end{figure}

\paragraph{The effect of pooling.} Before pooling, the benchmark scores at
ten results were 0.63, 0.85, 0.87, and 0.04. Pooling raised every engine and
changed no ranking. It mattered most where labels were thinnest: on real
questions LeanExplore rose from 0.33 to 0.50, LeanSearch from 0.58 to 0.77,
and Lean Finder from 0.62 to 0.80. Scores from single-answer labels understate
all engines, by different amounts.

\paragraph{The effect of weighting.} Weighted by the shares of
Table~\ref{tab:forms}, an earlier build of the written queries put the three
semantic engines level on Mathlib, at 0.82 to 0.83 at ten results. The forms
on which LeanExplore is weakest are the ones that its users rarely send it.
With equal buckets the same engines are twenty points apart. The choice of
weights decides the outcome, which is why we fix it by a rule that refers to
no engine.

%=====================================================================
\section{Limitations}
\label{sec:limitations}

\paragraph{Labels come from language models.} Queries, acceptance, and pooled
grades are all produced by language models, and no person has reviewed a
sample. The writer and the first judge are from different families, but the
pooled grades come from the family of the main writer. Language models are
reported to agree with human assessors well enough to rank systems, but less
well query by query~\cite{thomas2024llm,faggioli2023perspectives}. Graders also reported
original labels that look generous or wrong, which we have not corrected.

\paragraph{The pool covers four engines.} A new engine may return a valid
answer that none of the four did, and it counts as a miss until it is judged.
Pools are known to be robust to this for systems similar to those
pooled~\cite{zobel1998reliable,buckley2004incomplete}, and less so for systems
that are different in kind.

\paragraph{The writer sees LeanExplore's descriptions.} The description shown
to the writer is the one LeanExplore indexes, so written queries may share
wording with its index. This favours LeanExplore in the first three buckets.
The real questions bucket is free of it, and it is the bucket in which
LeanExplore does worst relative to the others.

\paragraph{The frame comes from one engine.} The eight forms and the lengths
of written queries come from LeanExplore's traffic. The equal weighting
removes the shares from the score but not the vocabulary of forms.

\paragraph{Hosted engines are not pinned.} Each service was queried as it
stood on the dates given and may change. We release the results so that the
scores can be recomputed, but a later run of the same engine may differ.

\paragraph{What the benchmark leaves out.} Every query has an answer, so it
does not measure whether an engine can tell that nothing matches. It covers
Mathlib alone, at one date. It also measures a single search and not the
success of a person or agent who searches several times.

%=====================================================================
\section{Related work}

\paragraph{Evaluations of Lean search engines.} LeanSearch was evaluated on
50 queries written by its authors~\cite{gao2024leansearch}, and its second
version on larger sets built for it~\cite{gao2026leansearchv2}. Lean Finder
was evaluated on informalized statements, synthetic user queries, perturbed
formal statements, and proof states, each with one gold declaration that must
be matched exactly~\cite{lu2025leanfinder}. LeanExplore was evaluated on 300
generated topical queries with no labelled answers, on which a language model
ranked the engines~\cite{asher2025leanexplore}. Each set was built with the
engine under test and fits its style of query. MIRB~\cite{ju2025mirb} collects
mathematical retrieval tasks, some of them formal. Premise selection
benchmarks~\cite{yang2023leandojo,tao2025premise} ask for the lemmas a proof
uses given its goal, which is the task of our \code{goal} form with labels
taken from proofs.

\paragraph{Retrieval benchmarks.} BEIR~\cite{thakur2021beir} and
MTEB~\cite{muennighoff2023mteb} showed that a ranking of retrieval models on
one task does not carry over to another, and report an unweighted mean over
tasks, as we do over buckets. CodeSearchNet~\cite{husain2019codesearchnet}
derives queries from documentation, and CoSQA~\cite{huang2021cosqa} pairs
real web queries with code. Our forms play the part of the query classes of
web search~\cite{broder2002taxonomy}, read from the surface of the query.

\paragraph{Test collections.} Pooling and the reuse of pooled judgments go
back to the proposal of \citet{sparckjones1975pooling} and to
TREC~\cite{voorhees2005trec}. Graded relevance follows
\citet{jarvelin2002cumulated}. Language-model queries and judgments produce
system rankings close to those of human
collections~\cite{rahmani2024synthetic,thomas2024llm}, with the cautions of
\citet{faggioli2023perspectives}.

%=====================================================================
\section{Conclusion}

A benchmark for Lean declaration search has to decide which queries count,
which declarations count, and which answers count, and each decision can
favour an engine. We weight four kinds of search equally, admit only
declarations that every engine indexes, and complete the answers by pooling.
On 2,560 test queries Lean Finder and LeanSearch are close, and both are well
ahead of LeanExplore, most of all on formal statements, proof states, and
questions from real discussions. Weighted by LeanExplore's own traffic the
three had looked level. Human review of a sample of labels and a track for
queries with no answer are the next steps.

{\small
\bibliographystyle{plainnat}
\bibliography{references}
}

\clearpage
\appendix

\section{Scores by form}
\label{app:forms}

Table~\ref{tab:byform} gives the numbers behind Figure~\ref{fig:forms}.

\begin{table}[H]
\centering\small
\caption{Success@10 on the test split by form, with Success@1 in brackets.}
\label{tab:byform}
\begin{tabular}{@{}lrrrrr@{}}
\toprule
\thead{Form} & \thead{Queries} & \thead{LeanExplore} & \thead{LeanSearch} & \thead{Lean Finder} & \thead{Loogle} \\
\midrule
\code{name} & 128 & 0.89 (0.65) & 0.85 (0.72) & 0.80 (0.59) & 0.31 (0.31) \\
\rowrule
\code{names} & 64 & 0.88 (0.31) & 0.95 (0.42) & 0.94 (0.56) & 0.11 (0.05) \\
\rowrule
\code{names+keywords} & 192 & 0.92 (0.45) & 0.94 (0.61) & 0.97 (0.82) & 0.00 (0.00) \\
\rowrule
\code{keywords} & 256 & 0.93 (0.56) & 0.96 (0.63) & 0.99 (0.82) & 0.00 (0.00) \\
\rowrule
\code{prose} & 384 & 0.86 (0.50) & 0.98 (0.68) & 0.99 (0.86) & 0.00 (0.00) \\
\rowrule
\code{prose+names} & 256 & 0.87 (0.45) & 0.97 (0.77) & 0.98 (0.91) & 0.00 (0.00) \\
\rowrule
\code{formal} & 384 & 0.59 (0.20) & 0.95 (0.70) & 0.97 (0.88) & 0.16 (0.16) \\
\rowrule
\code{goal} & 256 & 0.46 (0.06) & 0.95 (0.65) & 0.96 (0.90) & 0.00 (0.00) \\
\rowrule
\code{question} & 640 & 0.50 (0.27) & 0.77 (0.45) & 0.80 (0.56) & 0.00 (0.00) \\
\bottomrule
\end{tabular}
\end{table}

\section{Searches by source}
\label{app:sources}

Table~\ref{tab:sources} divides the shares of Table~\ref{tab:forms} by where
the search came from. The programmatic interface accounts for 77\% of
searches, the website for 20\%, and the Model Context Protocol server for
3\%.

\begin{table}[H]
\centering\small
\caption{Share of LeanExplore searches in each form, by source. A dash means
that fewer than ten searches were seen.}
\label{tab:sources}
\begin{tabular}{@{}lrrrrr@{}}
\toprule
\thead{Form} & \thead{All} & \thead{Website} & \thead{API} & \thead{MCP} & \thead{Median words} \\
\midrule
\code{name} & 5.5\% & 6.4\% & 4.2\% & 27.8\% & 1 \\
\rowrule
\code{names} & 1.5\% & 0.4\% & 1.7\% & 4.2\% & 2 \\
\rowrule
\code{names+keywords} & 22.7\% & 9.6\% & 26.4\% & 15.7\% & 5 \\
\rowrule
\code{keywords} & 26.2\% & 21.0\% & 28.0\% & 16.2\% & 6 \\
\rowrule
\code{prose} & 23.9\% & 46.2\% & 17.9\% & 26.3\% & 11 \\
\rowrule
\code{prose+names} & 17.0\% & 13.9\% & 18.2\% & 8.3\% & 10 \\
\rowrule
\code{formal} & 2.8\% & 2.6\% & 3.0\% & 1.6\% & 11 \\
\rowrule
\code{goal} & 0.4\% & -- & 0.5\% & -- & 12 \\
\bottomrule
\end{tabular}
\end{table}

\end{document}
